% Check criterion, deliverable, and the Decision Log question of Day 1.
% The criterion and the deliverable come from the Day 1 section of the
% syllabus and from Labs/day01/README.md.
\begin{frame}{Check criterion}
  The instructor checks this at the end of the lab:
  \begin{itemize}
    \item \cmark\ All five sensor tests of Part B pass.
    \item \cmark\ The report has the two lines of the build output, and the
          values of \code{memory\_report} for the two PSRAM settings.
    \item \cmark\ The notebook prints \code{Task 1: complete}.
    \item \cmark\ The report has the predictions of task 2. You wrote the
          predictions before the measurement.
    \item \cmark\ The table has six rows and four devices. Each ``does not
          fit'' has a reason.
    \item \cmark\ The Decision Log gives numbers and names one trade-off.
  \end{itemize}
  \vskip 0.5em
  \textbf{Hand in:} the file \code{report.md}. Show the notebook and the
  report to the instructor.
  \vskip 0.5em
  \keypoint{A result with no reason is not complete. ``Does not fit'' is the
    start of the answer. The budget that decides is the answer.}
\end{frame}

\begin{frame}{Decision Log}
  Write about 100 words. State one design decision, give your measured
  numbers, and name the trade-off.
  \vskip 0.4em
  \begin{exerciseblock}
    You must classify images of 96 $\times$ 96 pixels on the XIAOML Kit.
    \begin{itemize}
      \item Which of the three models do you select?
      \item In which data type: \code{float32} or \code{int8}?
      \item With PSRAM or with no PSRAM?
      \item Which budget decides?
    \end{itemize}
  \end{exerciseblock}
  \vskip 0.2em
  A good Decision Log:
  \begin{itemize}
    \item gives the model size and the peak activation memory from your
          table,
    \item gives the flash budget and the RAM budget that you measured on
          your board,
    \item names what you lose with your decision. Example: memory against
          accuracy.
  \end{itemize}
\end{frame}
